\chapter{引言}

教育对一个国家的方方面面都起到重要作用，它与一个国家的文化、人民的生计和世界的和平发展息息相关。中国的持续发展使得人民对教育的要求不断提高，推动教育改革对建设具有中国特色的教育制度有十分重要的意义。在过去，对高考成绩的追求使人们往往集中注意力于能够提高成绩的教学手段，即“填鸭式教学”，而忽略了由此引发的种种问题，例如学生的创新能力。成绩的提升不能以牺牲学生创造力为代价，因此，寻求有效教学手段以促进学生成绩提升的同时还应培养学生自主思考的能力，协调学生成绩提升和“死记硬背”之间的矛盾是当前教育发展进程中亟待解决的问题。

从古至今，提问便是一种重要的教学手段。孔子曾经说过:“不愤不启，不悱不发，举一隅不以三隅反，则不复也。” 孔子认为，教师应当在学生冥思苦想而不得解、思绪万千而无以诉的时候启发学生。而在这启发式教学的过程中，最常用的教学手段便是诘问，即提问。苏格拉底使用“产婆术”帮助人们获取知识。在与别人谈话的时候，苏格拉底会不断发问，以揭示对方话语中自相矛盾之处，从具体事例出发，不断深入，逐步反驳错误思想，从而走向理性认识。

近年来，课堂提问受到广泛关注。不论是美国学者克拉克所提出的课堂提问的19种功能，还是还是国内学者姚安娣所提出的提问的7种作用，都在告诉我们，课堂提问不只是简单的一问一答的过程，而是一种具有深远影响与深刻意义的教学行为，值得我们不断深入探讨。

在现有的文献中，关于课堂提问的研究大多集中在普通教育方面，但很少有人研究数学课上的提问，更少有人研究数学复习课中的课堂提问。本文将以数学课堂教学录像作为研究对象，采用录像分析法，通过反复观看录像，对复习课这种数学课型的课堂提问进行收集，运用布鲁姆学习分类法对问题进行分类，进一步分析在复习课这一数学课型当中课堂提问应遵循何种原则。在新课改的背景下研究课堂提问这一教学手段，可以扩展教育学界对提问教学的认识，推进与之相关的有效性提问、提问数量、教师理答方式等研究进程，为数学教学的课堂提问设计提供必要的参考依据。